% GENERATED FILE. Edit canonical lessons, then run npm run generate:books.
% canonical-source-hash: fnv1a64:ee4c399f654e1858
% canonical-lessons: ZH-C410-changtu, ZH-C410-chefei, ZH-C410-tongguo, ZH-C410-kaitong, ZH-C410-shanglu

\chapter{Long Trips}
\label{ch:zh-long-trips}
\hlchaptermodality{410}

\begin{chapteropening}
\textbf{By the end of this chapter:} \emph{I can talk about long-distance coaches and fares, say a train passes through, a line opens and we set off.}

Complete the last lesson of chapter 410: I can talk about long-distance coaches and fares, say a train passes through, a line opens and we set off.
\end{chapteropening}

\section[chángtú]{chángtú — long-distance}

\label{lesson:ZH-C410-changtu}



\begin{quote}
Before the new one: say the Mandarin for a ticket seller, then the Mandarin for to take a taxi.
\end{quote}

\subsection*{You'll want to know: chángtú}
\emph{chángtú} — \textquotedblleft{}long-distance\textquotedblright{}.

\begin{quote}
\emph{zuò chángtú chē qù hěn lèi} — Going by long-distance coach is tiring.
\end{quote}

Its characters come later; for now it is a word for the ear.

\begin{cousinweb}
\emph{chángtú} is built from \zh{长} 'long' and \emph{tú} 'a road', both inherited from Middle Chinese, \emph{drjang} and \emph{du} (Guangyun, in Baxter's transcription), which glosses the second 'a road'.
\end{cousinweb}

\begin{grammarlens}[title={What you've built}]
One more word for getting around.
\end{grammarlens}

\subsection*{Your turn}
\begin{itemize}
\raggedright
  \item \emph{Say it:} \emph{chángtú}
  \item \emph{Say it:} \emph{chángtú}, once more
  \item \emph{Say it:} \emph{zuò chángtú chē qù hěn lèi} — Going by long-distance coach is tiring.
  \item \emph{Recall:} say the Mandarin for a ticket seller, then the Mandarin for to take a taxi, then say \emph{chángtú} again
\end{itemize}

\begin{tcolorbox}[breakable,colback=teal!4,colframe=teal!35!black,title={Before you move on}]
What does chángtú mean? (\textquotedblleft{}Long-distance\textquotedblright{}.) Say it once more.
\end{tcolorbox}
\section[chēfèi]{chēfèi — fare for a ride}

\label{lesson:ZH-C410-chefei}



\begin{quote}
Before the new one: say the Mandarin for to take a taxi, then the Mandarin for long-distance.
\end{quote}

\subsection*{You'll want to know: chēfèi}
\emph{chēfèi} — \textquotedblleft{}fare for a ride\textquotedblright{}.

\begin{quote}
\emph{zuò chūzūchē qù jīchǎng chēfèi shì duōshao qián} — How much is the taxi fare to the airport?
\end{quote}

Its characters come later; for now it is a word for the ear.

\begin{cousinweb}
\emph{chēfèi} is built from \zh{车} 'vehicle' and \emph{fèi} 'a cost', both inherited from Middle Chinese, \emph{tsyhae} and \emph{phj+jH} (Guangyun, in Baxter's transcription), which glosses the second 'to use up'.
\end{cousinweb}

\begin{grammarlens}[title={What you've built}]
One more word for getting around.
\end{grammarlens}

\subsection*{Your turn}
\begin{itemize}
\raggedright
  \item \emph{Say it:} \emph{chēfèi}
  \item \emph{Say it:} \emph{chēfèi}, once more
  \item \emph{Say it:} \emph{zuò chūzūchē qù jīchǎng chēfèi shì duōshao qián} — How much is the taxi fare to the airport?
  \item \emph{Recall:} say the Mandarin for to take a taxi, then the Mandarin for long-distance, then say \emph{chēfèi} again
\end{itemize}

\begin{tcolorbox}[breakable,colback=teal!4,colframe=teal!35!black,title={Before you move on}]
What does chēfèi mean? (\textquotedblleft{}Fare for a ride\textquotedblright{}.) Say it once more.
\end{tcolorbox}
\section[\texorpdfstring{\zh{通过} (tōngguò)}{通过 (tōngguò)}]{\zh{通过} (tōngguò) — to pass through}

\label{lesson:ZH-C410-tongguo}



\begin{quote}
Before the new one: say the Mandarin for long-distance, then the Mandarin for the fare for a ride.
\end{quote}

\subsection*{You'll want to know: \zh{通过}}
\textbf{\zh{通过}} — \emph{tōngguò} — \textquotedblleft{}to pass through\textquotedblright{}.

\begin{quote}
\textbf{\zh{检票以后我们从这儿通过}} \emph{jiǎnpiào yǐhòu wǒmen cóng zhèr tōngguò} — After the ticket check we go through here.
\end{quote}

\begin{cousinweb}
\textbf{\zh{通过}} (\emph{tōngguò}) is built from \zh{通} 'to go through' and \zh{过} 'to cross', both inherited from Middle Chinese, \emph{thuwng} and \emph{kwaH} (Guangyun, in Baxter's transcription), which glosses the second 'to go beyond'.
\end{cousinweb}

\begin{grammarlens}[title={What you've built}]
One more word for getting around.
\end{grammarlens}

\subsection*{Your turn}
\begin{itemize}
\raggedright
  \item \emph{Say it:} \emph{tōngguò}
  \item \emph{Say it:} \emph{tōngguò}, once more
  \item \emph{Say it:} \emph{jiǎnpiào yǐhòu wǒmen cóng zhèr tōngguò} — After the ticket check we go through here.
  \item \emph{Recall:} say the Mandarin for long-distance, then the Mandarin for the fare for a ride, then say \emph{tōngguò} again
\end{itemize}

\begin{tcolorbox}[breakable,colback=teal!4,colframe=teal!35!black,title={Before you move on}]
What does \zh{通过} mean? (\textquotedblleft{}To pass through\textquotedblright{}.) Say it once more.
\end{tcolorbox}
\section[\texorpdfstring{\zh{开通} (kāitōng)}{开通 (kāitōng)}]{\zh{开通} (kāitōng) — to launch a new transport line}

\label{lesson:ZH-C410-kaitong}



\begin{quote}
Before the new one: say the Mandarin for the fare for a ride, then the Mandarin for to pass through.
\end{quote}

\subsection*{You'll want to know: \zh{开通}}
\textbf{\zh{开通}} — \emph{kāitōng} — \textquotedblleft{}to launch a new transport line\textquotedblright{}.

\begin{quote}
\emph{wǒmen zhèr kāitōng le gāotiě} — High-speed rail has opened in our area.
\end{quote}

\begin{cousinweb}
\textbf{\zh{开通}} (\emph{kāitōng}) is built from \zh{开} 'to open' and \zh{通} 'to go through', both inherited from Middle Chinese, \emph{khoj} and \emph{thuwng} (Guangyun, in Baxter's transcription).
\end{cousinweb}

\begin{grammarlens}[title={What you've built}]
One more word for getting around.
\end{grammarlens}

\subsection*{Your turn}
\begin{itemize}
\raggedright
  \item \emph{Say it:} \emph{kāitōng}
  \item \emph{Say it:} \emph{kāitōng}, once more
  \item \emph{Say it:} \emph{wǒmen zhèr kāitōng le gāotiě} — High-speed rail has opened in our area.
  \item \emph{Recall:} say the Mandarin for the fare for a ride, then the Mandarin for to pass through, then say \emph{kāitōng} again
\end{itemize}

\begin{tcolorbox}[breakable,colback=teal!4,colframe=teal!35!black,title={Before you move on}]
What does \zh{开通} mean? (\textquotedblleft{}To launch a new transport line\textquotedblright{}.) Say it once more.
\end{tcolorbox}
\section[\texorpdfstring{\zh{上路} (shànglù)}{上路 (shànglù)}]{\zh{上路} (shànglù) — to set out on a journey; to hit the road}

\label{lesson:ZH-C410-shanglu}



\begin{quote}
Before the new one: say the Mandarin for to pass through, then the Mandarin for to launch a new transport line.
\end{quote}

\subsection*{You'll want to know: \zh{上路}}
\textbf{\zh{上路}} — \emph{shànglù} — \textquotedblleft{}to set out on a journey; to hit the road\textquotedblright{}.

\begin{quote}
\textbf{\zh{我们早上六点上路}} \emph{wǒmen zǎoshang liù diǎn shànglù} — We set off at six in the morning.
\end{quote}

\begin{cousinweb}
\textbf{\zh{上路}} (\emph{shànglù}) is built from \zh{上} 'to go up' and \zh{路} 'a road', both inherited from Middle Chinese, \emph{dzyangX} and \emph{luH} (Guangyun, in Baxter's transcription), which glosses the first 'to climb'.
\end{cousinweb}

\begin{grammarlens}[title={What you've built}]
One more word for getting around.
\end{grammarlens}

\subsection*{Your turn}
\begin{itemize}
\raggedright
  \item \emph{Say it:} \emph{shànglù}
  \item \emph{Say it:} \emph{shànglù}, once more
  \item \emph{Say it:} \emph{wǒmen zǎoshang liù diǎn shànglù} — We set off at six in the morning.
  \item \emph{Recall:} say the Mandarin for to pass through, then the Mandarin for to launch a new transport line, then say \emph{shànglù} again
\end{itemize}

\begin{tcolorbox}[breakable,colback=teal!4,colframe=teal!35!black,title={Before you move on}]
What does \zh{上路} mean? (\textquotedblleft{}To set out on a journey; to hit the road\textquotedblright{}.) Say it once more.
\end{tcolorbox}
